\par\bigskip\Needspace{7\baselineskip}
\pdfbookmark[1]{Handbook of Navigation. 1850.}{schiffahrtskunde-1850}
\begin{center}\textit{Handbuch der Schiffahrtskunde} by C. \textsc{Rümker}. 1850. Page 76.\par\medskip
\textit{Solution of a geometrical problem.}\end{center}
\hspace*{\parindent}At three points \((1),(2),(3)\), lying on a straight line \(\mathrm{(I)}\) at known distances from one another, \(A\) from \((1)\) to \((2)\) and \(B\) from \((2)\) to \((3)\), the angles \(\vartheta,\vartheta',\vartheta''\) subtended by two other points \((4),(5)\), whose mutual distance \(=2c\) is also known, have been measured. The positions of the first three points relative to the last two are sought. To leave nothing undetermined, I assume that all three angles are measured with the same sense of increase from \((4)\) to \((5)\), that distances on line \(\mathrm{(I)}\) are counted as positive in a specified direction (so that if, for any reason, the point lying between the other two were given a label other than \((2)\), then \(A\) and \(B\) would have different signs), and that \(c\) is to be taken as positive.
